\subsection{Sensitivity filtering}
\label{subsec:sensitivity_filtering}

One of the most severe issues in topology optimization is the checkerboard pattern, which is a common artifact that appears in the optimized designs.
In order to mitigate this issue, it is common practice to apply a sensitivity filtering technique to the optimization algorithm.

In particular, a commonly used technique to reduce the risk of checkerboard patterns is the so-called sensitivity filtering, in which the sensitivities are convoluted with a filter kernel.
In other words, the density of each element is updated based on the average sensitivities of its neighboring elements, rather than the sensitivities of the element itself.

Such a filter, can be defined over the compliance sensitivities as:

\begin{equation}
    \frac{\hat{\partial C}}{\partial x_j} = \frac{1}{x_j \sum_{j} H_{j}} \sum_{j} H_{j} x_j \frac{\partial C}{\partial x_j}
    \label{eq:sensitivity_filtering}
\end{equation}

Where $H_{j}$ is the filter kernel, and the sum is performed over the neighboring elements of the element $j$.
The filter kernel $H_{j}$ can be defined in different ways, but a common choice is the following:

\begin{equation}
    H_{j} = r_{min} - dist(j, i) \quad \{i \in N \mid dist(j, i) \leq r_{min}\}
    \label{eq:filter_kernel}
\end{equation}
